% Generated draft: requires mathematical and source review.
\documentclass[11pt]{article}
\usepackage[margin=25mm]{geometry}
\usepackage{fontspec}
\setmainfont{Latin Modern Roman}
\usepackage{amsmath,amssymb,mathtools,unicode-math}
\setmathfont{Latin Modern Math}
\usepackage{enumitem,needspace,xurl,hyperref}
\setlength{\parindent}{0pt}
\setlength{\parskip}{5pt}
\setlength{\emergencystretch}{3em}
\newcommand{\N}{\mathbb{N}}
\newcommand{\Z}{\mathbb{Z}}
\newcommand{\Q}{\mathbb{Q}}
\newcommand{\R}{\mathbb{R}}
\newcommand{\C}{\mathbb{C}}
\newcommand{\F}{\mathbb{F}}
\newcommand{\A}{\mathbb{A}}
\newcommand{\PP}{\mathbb{P}}
\newcommand{\E}{\mathbb{E}}
\newcommand{\eps}{\varepsilon}
\newcommand{\dd}{\,\mathrm{d}}
\begin{document}
\section*{UnsolvedMath ID: 4600013}
\textbf{Source problem code:} catalogue entry
\subsection{Definitions and mathematical statement}
Let $\alpha:\mathbb Z^2\to\operatorname{Homeo}(X)$ be an action on a compact metric space $(X,d)$. Assume that for some nonzero $\boldsymbol n\in\mathbb Z^2$, the homeomorphism $\alpha^{\boldsymbol n}$ is topologically conjugate to a two-sided shift of finite type, that is, a shift on sequences over a finite alphabet determined by finitely many forbidden words. No mixing assumption is stated.

A line $L\subset\mathbb R^2$ through the origin is expansive for $\alpha$ if there exist $r>0$ and $\varepsilon>0$ such that, whenever $x,y\in X$ satisfy
\[
d(\alpha^{\boldsymbol m}x,\alpha^{\boldsymbol m}y)\leq\varepsilon
\quad\text{for all }\boldsymbol m\in\mathbb Z^2
\text{ with }\operatorname{dist}(\boldsymbol m,L)\leq r,
\]
then $x=y$. Let
\[
\mathcal E(\alpha)=\{v\in\mathbb R^2\setminus\{0\}:\mathbb Rv\text{ is expansive}\}.
\]
Its connected components are the expansive components in the cone convention. Equivalently, one may describe their directional intervals; identifying opposite directions does not change either existence question below. A line has irrational slope precisely when it contains no nonzero vector of $\mathbb Z^2$.

Under the stated shift-of-finite-type-element hypothesis:
\begin{enumerate}
\item Can $\mathcal E(\alpha)$ have infinitely many connected components?
\item Can an expansive component have a nonzero boundary ray contained in a line of irrational slope?
\end{enumerate}
The supplied research summary describes both questions as unresolved through its 2026 assessment. Its phrase ``$\mathbb Z^2$ SFTs'' is not used to strengthen the actual hypothesis, which says that one element of the action is a one-dimensional SFT.
\subsection{Short English statement}
For a $\mathbb Z^2$ action with a shift-of-finite-type element, can the set of expansive directions have infinitely many components? Can one such component end at an irrational direction?
\subsection{Sources}
\begin{itemize}
\item Boyle, \emph{Open Problems in Symbolic Dynamics} (2008), item 11.2, \url{https://www.math.umd.edu/~mboyle/papers/openfinalsub3nov2007.pdf}; accessed 2026-07-29, with PDF-text extraction recorded in the supplied provenance.
\item The supplied research summary attributes the expansive-component framework to Boyle--Lind, but supplies no separate publication identifier for that work.
\item Local registry \texttt{data/unsolvedmath/problems.json}, permanent UnsolvedMath ID $4600013$, source code catalogue entry, \emph{Expansive components}. Dataset attribution: Ulam AI and the UnsolvedMath Contributors.
\end{itemize}
\end{document}
